% ---------------------------------------------------------------------------
% Shared preamble for the lecture notes
% ---------------------------------------------------------------------------
%
% NOTE: main.tex uses the `lecturenotes` document class (notes/lecturenotes.cls),
% which already loads hyperref, amsmath (intlimits), amssymb, amsthm, tikz,
% mathtools, babel, etc., and already defines the theorem-like environments
% theorem/lemma/proposition/corollary/definition/example/remark (plus
% exercise/claim/conjecture/fact/problem/notation). Do not reload or redefine
% those here. The class requires LuaLaTeX and does not support inputenc/fontenc.

\usepackage{graphicx}
\usepackage{cleveref}

% Bibliography: biblatex + biber. Needed because classic BibTeX resolves .bib paths
% relative to the compiler's working directory (unreliable on Overleaf), while biblatex
% resolves \addbibresource relative to the main .tex file itself — which is what lets us
% keep bibliography.bib at the repo root and reference it as ../bibliography.bib from here.
\usepackage[backend=biber, style=numeric, sorting=none]{biblatex}
\addbibresource{../bibliography.bib}

% The class sets a narrow text block (right=2in, to leave room for \margintext),
% which leaves TeX little slack on lines carrying long unbreakable tokens --
% arXiv identifiers and DOIs in the bibliography, mostly. A little emergency
% stretch avoids overfull lines there without loosening normal body text.
\setlength{\emergencystretch}{2em}

\usetikzlibrary{shapes,positioning,arrows.meta}

% --- Cross-references to the four course parts ---
% The parts are \section's in a document whose section 1 is "General
% Information", so \ref{sec:partN} would print the section number (2,3,4,5)
% rather than the part number. \partref{n} prints the Roman numeral and still
% hyperlinks: "Part~\partref{1}", "Parts~\partref{1} and~\partref{2}".
\newcommand{\partref}[1]{%
  \ifcase#1\or
    \hyperref[sec:part1]{I}\or
    \hyperref[sec:part2]{II}\or
    \hyperref[sec:part3]{III}\or
    \hyperref[sec:part4]{IV}%
  \fi}

% --- Tensor-network notation macros (placeholders — adapt to your convention) ---
\newcommand{\tensor}[1]{\mathsf{#1}}
\newcommand{\rank}{\chi}
\newcommand{\bond}{r}
